\begin{wrapfigure}{l}{0.1\textwidth}
	\vspace{-10pt}
	\includegraphics[width=\linewidth]{images/ikony_ubuntutweak.png}
\end{wrapfigure}

Program \textcolor{ubuntu_orange}{Ubuntu Tweak} jest zestawem narzędzi pozwalających usprawnić system Ubuntu. W przeciwieństwie do omawianego wcześniej Unity Tweak Tool, ten program mniej koncentruje się na wyglądzie systemu, a bardziej na jego zachowaniu. Część opcji Ubuntu Tweak pokrywa się z Unity Tweak Tool (głównie tych dotyczących wyglądu środowiska), ale ma też kilka bardzo przydatnych narzędzi:

\begin{itemize}
\item \menu{{Usprawnienia}>{Logowanie}} zmiana wyglądu ekranu logowania (Greetera);
\item \menu{{Zarządzanie}>{Polecenia aktywatorów}} edycja Szybkich List (QuickList) aktywatorów na pasku Launchera;
\item \menu{{Zarządzanie}>{Menadżer typów plików}} zarządzanie powiązaniami plików z odpowiednimi programami;
\item \menu{Czyszczenie} podstawowe sprzątanie systemu i przestrzeni użytownika.
\end{itemize}

Ubuntu Tweak możesz pobrać bezpośrednio ze \href{http://ubuntu-tweak.com/}{strony autora} lub dodać odpowiednie repozytorium i zainstalować przy pomocy menadżera pakietów:
\begin{lstlisting}[language=bash]
sudo add-apt-repository ppa:tualatrix/ppa
sudo apt-get update
sudo apt-get install ubuntu-tweak
\end{lstlisting}

Więcej o instalacji oprogramowania w Ubuntu przeczytasz w rozdziale \ref{instalacja_oprogramowania}: ,,Instalacja oprogramowania''.
